\documentclass[10pt,a4paper]{amsart}
\usepackage[margin=19mm]{geometry}
\usepackage{amsmath,amssymb,mathtools}
\usepackage[scr=rsfs,cal=euler]{mathalfa}
\usepackage{xfrac}
\usepackage[unicode,hidelinks,bookmarks=false]{hyperref}
\newcommand{\ie}{i.e.\ }
\newcommand{\cf}{cf.\ }
\newcommand{\resp}{respectively\ }
\newcommand{\Subsection}{\subsection}
\providecommand{\nobreakdash}{\nobreak\hbox{-}}
\setlength{\emergencystretch}{2em}
\multlinegap0pt
\usepackage{bm}
\usepackage{bbm}
\usepackage{dsfont}
\usepackage{xspace}
\usepackage{cancel}
\usepackage{mathtools}
\usepackage{amsthm}
\makeatletter
\@ifundefined{theorem}{\newtheorem{theorem}{Theorem}}{}
\@ifundefined{lemma}{\newtheorem{lemma}[theorem]{Lemma}}{}
\makeatother
\def\D{\mathbb{B}^n_+}
\def\bn{\mathbb{B}^n}
\def\cC{\mathcal{C}}
\def\cN{\mathcal{N}}
\def\dist{\mathrm{dist}}
\def\eps{\varepsilon}
\def\hn{\mathbb{R}^n_+}
\def\idhn{\int_{\partial\hn}}
\def\ihn{\int_{\hn}}
\def\irn{\int_{\rn}}
\def\n{\mathbb{N}}
\def\rh{\rightharpoonup}
\def\rn{\mathbb{R}^n}
\def\r{\mathbb{R}}
\def\supp{\mathrm{supp}}
\def\vp{\varphi}
\def\what{\widehat}
\title[Sign-changing solutions to the Yamabe problem on a spherical]{Sign-changing solutions to the Yamabe problem on a spherical cap}
\author{M\'onica Clapp, Benedetta Pellacci, Angela Pistoia}
\date{Source: arXiv:2605.29045v1 (CC BY 4.0)}
\begin{document}
\maketitle

\noindent\emph{Source and licence.} Sign-changing solutions to the Yamabe problem on a spherical cap. Version arXiv:2605.29045v1 (CC BY 4.0), distributed under the Creative Commons Attribution 4.0 International licence, as recorded in the arXiv licence metadata for this exact version and shown on the abstract page https://arxiv.org/abs/2605.29045v1. Full rights notice: licenses/arxiv-2605.29045-CC-BY-4.0.txt.\par\smallskip

\noindent\textit{Editorial scope.} Counted unit: the Theorem whose source label is \texttt{thm:minimizing} in the article (source0.tex lines 392--403, source label \texttt{thm:minimizing}), delivered together with the article's own complete original proof of it (source0.tex lines 405--588, one \texttt{proof} environment of 184 lines). No mathematics is written, repaired or re-derived: statement and proof are the article's own text line for line. Required context retained uncounted: problem\_cap (lines 162--167); eq:norm (lines 230--232); eq:F (lines 270--272); lem:nehari (lines 274--280); problem\_rn (lines 289--294); lem:no\_minimum (lines 356--358); lem:G (lines 368--374). Every definition, display and label of those lines is retained unchanged. Omitted from this package and not redistributed: 1--161, 168--229, 233--269, 273--273, 281--288, 295--355, 359--367, 375--391, 404--404, 589--670. Nothing inside a retained excerpt is omitted or altered: every retained body is delivered exactly as the article's lines are written, so no omission receipt of any kind accompanies this package. Citations: the retained excerpts carry no citation command at all, so no bibliography entry of the article is transcribed and no reference is redistributed. Reference adaptation: none was needed and none was made; every reference inside the delivered unit resolves inside it, so no unresolved reference or citation is produced. Printed slips kept literal: no printed word, symbol, hypothesis or number of the counted statement or of its proof was altered. Layout and class: the article's own class is \texttt{amsart} with packages including inputenc, amsmath, amsfonts, amssymb, graphicx, verbatim, mathrsfs, upref, amsthm, amsxtra, exscale, cite, dsfont, geometry; the wrapper is the lane's pinned \texttt{amsart} editorial wrapper at 10pt on A4, which restates the theorem-like environments the retained text needs and adds the article's own macro definitions that the retained text needs (\texttt{\textbackslash D}, \texttt{\textbackslash bn}, \texttt{\textbackslash cC}, \texttt{\textbackslash cN}, \texttt{\textbackslash dist}, \texttt{\textbackslash eps}, \texttt{\textbackslash hn}, \texttt{\textbackslash idhn}, \texttt{\textbackslash ihn}, \texttt{\textbackslash irn}, \texttt{\textbackslash n}, \texttt{\textbackslash r}, \texttt{\textbackslash rh}, \texttt{\textbackslash rn}, \texttt{\textbackslash supp}, \texttt{\textbackslash vp}, \texttt{\textbackslash what}), with extra wrapper packages (bm, bbm, dsfont, xspace, cancel, mathtools). The delivered numbering is the unit's own. Assets: no retained span contains a figure environment, a graphics-inclusion command, a PDF-page inclusion, a table or a TikZ picture; no image file is copied into this package, the package declares no asset dependency and no image file is redistributed with it. Licence: version arXiv:2605.29045v1 of this article is distributed under the Creative Commons Attribution 4.0 International licence, as recorded in the arXiv licence metadata for this exact version and shown on the abstract page https://arxiv.org/abs/2605.29045v1. A full rights notice is delivered at licenses/arxiv-2605.29045-CC-BY-4.0.txt. Characters: every retained excerpt is pure ASCII, so the delivered engine, XeLaTeX with its default Latin Modern text font, reports no missing character.
\begin{equation}\label{problem_cap}
\begin{cases}
 -\Delta u = a_n|u|^\frac{4}{n-2}u &\text{in \ } \hn,\\
 \frac{\partial u}{\partial x_n} = -b|u|^\frac{2}{n-2}u &\text{on \ }\partial\hn,
\end{cases}
\end{equation}

\begin{equation}\label{eq:norm}
\|u\|:=\Big(\ihn|\nabla u|^2\Big)^\frac{1}{2}.
\end{equation}

\begin{equation}\label{eq:F}
F_b(u):=a_n\ihn |u|^p + b\idhn|u|^q.
\end{equation}

\begin{lemma}\label{lem:nehari}
\begin{itemize}
\item[$(a)$]There exists $c_0>0$ such that $\|u\|\geq c_0$ for all $u\in\cN^\phi_b(\hn)$.
\item[$(b)$]$\cN^\phi_b(\hn)$ is a Hilbert submanifold of class $\cC^2$ of $D^{1,2}(\hn)^\phi$ and a natural constraint for $J_b$.
\item[$(c)$]If $u\in D^{1,2}(\hn)\smallsetminus\{0\}$, then there exists a unique $t_u\in(0,\infty)$ such that $t_uu\in\cN^\phi_b(\hn)$. The function $t\mapsto J(tu)$ is strictly increasing in $[0,t_u]$ and strictly decreasing in $[t_u,\infty)$.
\end{itemize}
\end{lemma}

\begin{equation}\label{problem_rn}
\begin{cases}
 -\Delta u = a_n|u|^\frac{4}{n-2}u,\\
 u\in D^{1,2}(\rn)^\phi,
\end{cases}
\end{equation}

\begin{lemma} \label{lem:no_minimum}
$\mu^\phi_b(\D)=\mu^\phi_b(\hn)$ and $\mu^\phi_b(\D)$ is not attained by $J_b$ on $\cN_b^\phi(\D)$.
\end{lemma}

\begin{lemma}\label{lem:G}
If $G$ satisfies $(G_1)$ then, any given sequences $(\eps_k)$ in $(0,\infty)$ and $(y_k)$ in $\r^{n-1}$ contain subsequences such that one of the following statements holds:
\begin{enumerate}
\item[$(i)$] either there exists $C_0>0$ such that $\eps_{k}^{-1}|y_k|<C_0$ for all $k\in\n$,
\item[$(ii)$] or, for each $m\in\n$, there exist $g_1,\ldots,g_m\in G$ such that $\eps_{k}^{-1}|g_iy_k-g_jy_k|\to\infty$ as $k\to\infty$ for any $i\neq j$, $1\leq i,j\leq m$.
\end{enumerate}
\end{lemma}

\begin{theorem} \label{thm:minimizing}
Let $(u_k)$ be a sequence in $\cN_b^\phi(\D)$ such that $J_b(u_k)\to \mu^\phi_b(\D)$. Then, after passing to a subsequence, one of the following statements holds true:
\begin{itemize}
\item[$(I)$] There exists a sequence of positive numbers $(\eps_k)$ and a nontrivial solution $w\in D^{1,2}(\hn)^\phi$ to the problem \eqref{problem_cap} such that
$$\lim_{k\to\infty}\left\|u_k-\eps_k^\frac{2-n}{2}w\left(\frac{\cdot}{\eps_k}\right)\right\| = 0$$
and $J_b(w)=\mu^\phi_b(\hn)=\mu^\phi_b(\D)$.

\item[$(II)$] There is a sequence of positive numbers $(\eps_k)$, a sequence of points $\xi_k=(0,t_k)$ in $\D$ and a nontrivial solution $w\in D^{1,2}(\rn)^\phi$ to the problem \eqref{problem_rn} such that
$$\eps_k^{-1}\dist(\xi_k,\partial\D)\to\infty,\qquad\lim_{k\to\infty}\left\|u_k-\eps_k^\frac{2-n}{2}w\left(\frac{\cdot - \xi_k}{\eps_k}\right)\right\| = 0$$
and $\mu^\phi_\infty\leq J_\infty(w)=\mu^\phi_b(\D)$.
\end{itemize}
\end{theorem}

\begin{proof}
In what follows, we consider $V(\D)\subset D^{1,2}(\hn)$ by means of trivial extension. 

Since $u_k\in\cN_b^\phi(\D)$ we have that
$$J_b(u_k)=\frac{1}{2}\ihn|\nabla u_k|^2 - \frac{a_n}{p}\ihn|u_k|^p - \frac{b}{q}\idhn |u_k|^q\geq\frac{1}{2}\|u_k\|^2 - \frac{1}{q}F_b(u_k)=\frac{q-2}{2q}\|u_k\|^2,$$
with $\|\cdot\|$ and $F_b$ as defined in \eqref{eq:norm} and \eqref{eq:F}. Since $J_b(u_k)\to \mu^\phi_b(\D)$, the sequence $(u_k)$ is bounded and, passing to a subsequence, $u_k\rh u$ weakly in $V(\D)^\phi$. Using Ekeland's variational principle, we may assume that $J'_b(u_k)\to 0$ in $V(\D)'$. Then, for every $\vp\in \cC^\infty_c(\bn)$,
$$o(1)=J'_b(u_k)\vp=\int_{\D}\nabla u_k\cdot\nabla \vp - a_n\int_{\D} |u_k|^{p-2}u_k\vp - b\int_{\partial\D} |u_k|^{q-2}u_k\vp,$$
and, passing to the limit, yields $J_b'(u)=0$ in $V(\D)'$. We claim that $u=0$. Indeed, using the Brezis-Lieb lemma we see that
\begin{align}
J_b(u_k)&=J_b(u_k-u)+J_b(u)+o(1), \label{eq:D1}\\
J'_b(u_k)u_k&=J'_b(u_k-u)[u_k-u]+J'_b(u)u+o(1)=J'_b(u_k-u)[u_k-u]+o(1),\nonumber
\end{align}
and, as a consequence,
\begin{equation}\label{eq:D2}
J_b(u_k-u)\geq\frac{1}{2}\|u_k-u\|^2-\frac{1}{q}F_b(u_k-u)=\frac{q-2}{2q}\|u_k-u\|^2+o(1).
\end{equation}
Equations \eqref{eq:D1} and \eqref{eq:D2} imply that 
\begin{equation}\label{eq:D3}
J_b(u_k)\geq J_b(u)+o(1).
\end{equation}
So, if $u\neq 0$, we would have that $u\in\cN_b^\phi(\D)$ and $J_b(u)=\mu_1(\D)$, contradicting Lemma \ref{lem:no_minimum}. This shows that $u_k\rh 0$ weakly in $V(\D)$. 

Consider the concentration function
$$Q_k(r):=\sup _{x\in\rn}\Big\{\int_{B_r(x)\cap\D}|u_k|^p+\int_{B_r(x)\cap\Gamma_1}|u_k|^q\Big\},$$
where $B_r(x)$ is the open ball of radius $r$ centered at $x$ in $\rn$. By Lemma \ref{lem:nehari}$(a)$, 
$$0<c_0\leq\|u_k\|^2=F_b(u_k)+o(1)\leq a_n\int_{\D}|u_k|^p+\max\{b,0\}\int_{\Gamma_1}|u_k|^q+o(1).$$ 
Hence, there exists $\kappa_0>0$ such that, after passing to a subsequence, 
\begin{equation}\label{eq:kappa0}
0<\kappa_0<\int_{\D}|u_k|^p+\int_{\Gamma_1}|u_k|^q\qquad\text{for all \ }k\in\n.
\end{equation}
Let $\kappa\in (0,\min\{1,\kappa_0\})$, to be fixed later. Then, there exist $\eps_k>0$ and $x_k\in\rn$ such that
\begin{equation} \label{eq:kappa}
\kappa=Q_k(\eps_k)=\int_{B_{\eps_k}(x_k)\cap\D}|u_k|^p+\int_{B_{\eps_k}(x_k)\cap\Gamma_1}|u_k|^q.
\end{equation}
We claim that $\eps_k\to 0$. 

To prove this claim, we argue by contradiction. Assume there is $0<\what{\eps}\leq\eps_k$ for all $k$ and let $\vp\in\cC^\infty_c(B_{\what{\eps}}(z))$ for some $z\in\rn$. Since $u_k\to0$ strongly in $L^2(\D)$, using \eqref{eq:kappa} we obtain
\begin{align} \label{eq:eps_to_0}
&\int_{\D}|\nabla(\vp u_k)|^2=\int_{\D}\nabla u_k\cdot\nabla(\vp^2 u_k) + o(1) \nonumber \\
&=a_n\int_{\D}\vp^2|u_k|^p+b\int_{\Gamma_1}\vp^2|u_k|^q+o(1)  \nonumber \\
&\leq C_{n,b}\Big(\int_{\D}|u_k|^{p-2}|\vp u_k|^2+\int_{\Gamma_1}|u_k|^{q-2}|\vp u_k|^2\Big)+o(1) \nonumber  \\
&\leq C_{n,b}\left[\Big(\int_{\D\cap B_{\what{\eps}}(z)}|u_k|^p\Big)^\frac{p-2}{p}\Big(\int_{\D}|\vp u_k|^p\Big)^\frac{2}{p} + \Big(\int_{\Gamma_1\cap B_{\what{\eps}}(z)}|u_k|^q\Big)^\frac{q-2}{q}\Big(\int_{\Gamma_1}|\vp u_k|^q\Big)^\frac{2}{q}\right] + o(1) \nonumber  \\
&\leq \overline{C}_{n,b} \Big[\kappa^\frac{2}{n} + \kappa^\frac{1}{n-1}\Big]\int_{\D}|\nabla(\vp u_k)|^2 + o(1).
\end{align}
So choosing $\kappa$ sufficiently small we get that $\int_{\D}|\nabla(\vp u_k)|^2 = o(1)$ and, hence, that
$$\int_{\D}|\vp u_k|^p+\int_{\Gamma_1}|\vp u_k|^q=o(1)\qquad\text{for all \ }\vp\in\cC^\infty_c(B_{\what{\eps}}(z)), \ z\in\rn.$$
Taking $z_1,\ldots,z_\ell\in\rn$ such that $\D\subset B_\frac{\what\eps}{2}(z_1)\cup\cdots\cup B_\frac{\what\eps}{2}(z_\ell)$ and $\vp_i\in\cC^\infty_c(B_{\what{\eps}}(z_i))$ such that $\vp_i(z)=1$ if $|z-z_i|\leq\frac{\what\eps}{2}$ yields
$$\int_{\D}|u_k|^p+\int_{\Gamma_1}|u_k|^q\leq\sum_{i=1}^\ell\Big(\int_{\D}|\vp_iu_k|^p+\int_{\Gamma_1}|\vp_iu_k|^q\Big)=o(1),$$
which contradicts \eqref{eq:kappa0}. This shows that $\eps_k\to 0$.

Our next task is to replace the points $x_k$ with more convenient concentration points $\xi_k$. To this end, first, we 
write $x_k=(y_k,t_k)\in\r^{n-1}\times\r$ and recall that Lemma \ref{lem:G} gives two alternatives. We claim that the alternative $(ii)$ is impossible. Indeed, arguing by contradiction, assume that, for each $m\in\n$, there exist $g_1,\ldots,g_m\in G$ such that $\eps_{k}^{-1}|g_iy_k-g_jy_k|\to\infty$ as $k\to\infty$ for any $i\neq j$. Then, for $k$ large enough, $B_{\eps_k}(g_ix_k)\cap B_{\eps_k}(g_jx_k)=\emptyset$ and, as a consequence,
\begin{equation*}
m\kappa=\sum_{i=1}^m\Big(\int_{B_{\eps_k}(g_ix_k)\cap\D}|u_k|^p+\int_{B_{\eps_k}(g_ix_k)\cap\Gamma_1}|u_k|^q\Big)\leq\int_{\D}|u_k|^p+\int_{\Gamma_1}|u_k|^q
\end{equation*}
which is impossible because $(u_k)$ is bounded in $V(\D)$. Hence, the alternative $(i)$ must hold, that is, after passing to a subsequence, there exists $C_0>0$ such that $\eps_k^{-1}|y_k|<C_0$ for all $k\in\n$. Set $\zeta_k=(0,t_k)\in\r^{n-1}\times\r$. Then, \eqref{eq:kappa} yields
\begin{equation} \label{eq:kappa1}
\kappa\leq \int_{B_{(C_0+1)\eps_k}(\zeta_k)\cap\D}|u_k|^p+\int_{B_{(C_0+1)\eps_k}(\zeta_k)\cap\Gamma_1}|u_k|^q.
\end{equation}
Now we claim that, after passing to a subsequence, there exist points $\xi_k$ and $C_1>1$ such that
\begin{equation} \label{eq:kappa2}
\kappa\leq \int_{B_{C_1\eps_k}(\xi_k)\cap\D}|u_k|^p+\int_{B_{C_1\eps_k}(\xi_k)\cap\Gamma_1}|u_k|^q
\end{equation}
and one of the following statements holds true:
\begin{enumerate}
\item $\xi_k=(0,0)$ for all $k\in\n$.
\item $\xi_k=(0,t_k)\in\D$ for all $k\in\n$, and $\eps_k^{-1}\dist(\xi_k,\partial\D)\to\infty$.
\item $\xi_k=(0,1)$ for all $k\in\n$.
\end{enumerate}
Indeed, after passing to a subsequence there are three possibilities: If $(\eps_k^{-1}|\zeta_k|)$ is bounded, we set $\xi_k:=(0,0)$. Then, \eqref{eq:kappa2} follows from \eqref{eq:kappa1}. If $\eps_k^{-1}\dist(\zeta_k,\partial\D)\to\infty$ we set $\xi_k:=\zeta_k$. Since $\dist(\zeta_k,\D)\leq(C_0+1)\eps_k$ we have that $\zeta_k\in\D$ and \eqref{eq:kappa2} is simply \eqref{eq:kappa1}. Finally, if $(\eps_k^{-1}|\zeta_k|)$ is unbounded and $(\eps_k^{-1}|\zeta_k-(0,1)|)$ is bounded we set $\xi_k:=(0,1)$ and \eqref{eq:kappa2} follows from \eqref{eq:kappa1}.

Set
$$D_k:=\eps_k^{-1}(\D-\xi_k)\qquad\text{and}\qquad\what\Gamma_k:=\eps_k^{-1}(\Gamma_1-\xi_k),$$
and define
$$w_k(z)=\eps_k^\frac{n-2}{2}u_k(\eps_kz+\xi_k)\qquad\text{if \ }z\in D_k.$$
Since $\xi_k$ is a fixed-point of $G$, the function $w_k$ is $\phi$-equivariant, and it follows from \eqref{eq:kappa} and \eqref{eq:kappa2} that
\begin{align} \label{eq:kappa3}
\kappa =\sup_{x\in\rn}\Big\{\int_{B_1(z)\cap D_k}|w_k|^p+\int_{B_1(z)\cap\what\Gamma_k}|w_k|^q\Big\}\leq \int_{B_{C_1}(0)\cap D_k}|w_k|^p+\int_{B_{C_1}(0)\cap\what\Gamma_k}|w_k|^q.
\end{align}
Furthermore, for any $\vp\in\cC^\infty_c(\rn)$, setting $\vp_k(z):=\vp\big(\frac{x-\xi_k}{\eps_k}\big)$ we see that
\begin{equation}\label{eq:derivative1}
\int_{D_k}\nabla w_k\cdot\nabla(\vp^2w_k)-a_n\int_{D_k}\vp^2|w_k|^p-b\int_{\what\Gamma_k}\vp^2|w_k|^q = J_b'(u_k)[\vp_k^2u_k]=o(1).
\end{equation}
Next, we analize the behavior of $(w_k)$ in each of the three cases stated above.

\underline{Case $(1)$}: $\xi_k=(0,0)$ for all $k\in\n$.

In this case, $\bigcup_{k\geq 1}D_k=\hn$ and $w_k\in D^{1,2}(\hn)^\phi$. Since the sequence $(w_k)$ is bounded in $D^{1,2}(\hn)$, a subsequence satisfies $w_k\rh w$ weakly in $D^{1,2}(\hn)$, $w_k\to w$ a.e. in $\hn$, $w_k\to w$ in $L^2_{loc}(\hn)$, and $w$ is $\phi$-equivariant. We claim that $w\neq 0$. Otherwise, we get a contradiction as follows: Let $\vp\in\cC^\infty_c(B_1(z))$ for some $z\in\rn$. Note that $\supp(\vp)\cap\hn\subset D_k$ for every large enough $k$. So following the argument in \eqref{eq:eps_to_0}, this time using \eqref{eq:kappa3} and \eqref{eq:derivative1}, we conclude that 
$$\ihn|\vp w_k|^p+\idhn|\vp w_k|^q = o(1)\qquad\text{for every \ }\vp\in\cC^\infty_c(B_1(z)), \ z\in\rn.$$
Taking $z_1,\ldots,z_\ell\in\rn$ such that $B_{C_1}(0)\subset B_\frac{1}{2}(z_1)\cup\cdots\cup B_\frac{1}{2}(z_\ell)$ and $\vp_i\in\cC^\infty_c(B_1(z_i))$ such that $\vp_i(z)=1$ if $|z-z_i|\leq\frac{1}{2}$ we obtain
$$\int_{B_{C_1}(0)\cap \hn}|w_k|^p+\int_{B_{C_1}(0)\cap\partial\hn}|w_k|^q\leq\sum_{i=1}^\ell\Big(\ihn|\vp_iw_k|^p+\idhn|\vp_iw_k|^q\Big)=o(1),$$
which contradicts \eqref{eq:kappa3}. This shows that $w\neq 0$.

Furthermore, if $\vp\in\cC^\infty_c(\rn)$, then $\supp(\vp)\cap\hn\subset D_k$ for all large enough $k$, and setting $\vp_k(z):=\vp\big(\frac{x-\xi_k}{\eps_k}\big)$ we see that
\begin{equation*}
\ihn\nabla w_k\cdot\nabla\vp-a_n\ihn|w_k|^{p-2}w_k\vp-b\idhn|w_k|^{q-2}w_k\vp = J_b'(u_k)[\vp_k|_{\D}]=o(1).
\end{equation*}
Passing to the limit yields $J_b'(w)=0$ in $D^{1,2}(\hn)'$, that is, $w$ is a nontrivial $\phi$-equivariant solution to \eqref{problem_cap}. Now, using the Brezis-Lieb lemma and arguing  as we did at the beginning of the proof, we see that
\begin{align}\label{eq:H1}
\mu^\phi_b(\D)+o(1)&=J_b(w_k)=J_b(w_k-w)+J_b(w)+o(1)\geq J_b(w_k-w)+\mu^\phi_b(\hn)+o(1), \\
J_b(w_k-w)&\geq\frac{1}{2}\|w_k-w\|^2-\frac{1}{q}F_b(w_k-w)=\frac{q-2}{2q}\|w_k-w\|^2+o(1),
\end{align}
and, using Lemma \ref{lem:no_minimum}, we obtain  
$$\lim_{k\to\infty}\left\|u_k-\eps_k^\frac{2-n}{2}w\left(\frac{\cdot}{\eps_k}\right)\right\|=\lim_{k\to\infty}\|w_k-w\| = 0$$
and $J_b(w)=\mu^\phi_b(\D)=\mu^\phi_b(\hn)$. Hence, in this case, statement $(I)$ holds true.

\underline{Case $(2)$}: $\xi_k=(0,t_k)\in\D$ for all $k\in\n$, and $\eps_k^{-1}\dist(\xi_k,\partial\D)\to\infty$.

In this case, $\bigcup_{k\geq 1}D_k=\rn$. Consider the extension operator $D^{1,2}(\hn)\to D^{1,2}(\rn)$,  $v\mapsto v^*$, where
\begin{equation*}
v^*(x',t):=
\begin{cases}
v(x',t) &\text{if \ }t>0,\\
v(x',-t) &\text{if \ }t<0,
\end{cases}
\end{equation*}
and define $\what w_k\in D^{1,2}(\rn)$ by
$$\what w_k(z):=\eps_k^\frac{n-2}{2}u_k^*(\eps_kz+\xi_k).$$
Note that $\what w_k$ is $\phi$-equivariant and $\what w_k(z)=w_k(z)$ if $z\in D_k$. Since $(\what w_k)$ is bounded in $D^{1,2}(\rn)$, a subsequence satisfies $\what w_k\rh w$ weakly in $D^{1,2}(\rn)$, $\what w_k\to w$ a.e. in $\rn$, $\what w_k\to w$ in $L^2_{loc}(\rn)$, and $w$ is $\phi$-equivariant. If $\vp\in\cC^\infty_c(B_1(z))$ for some $z\in\rn$, then $\supp(\vp)\subset D_k$ for every large enough $k$ and following the argument in \eqref{eq:eps_to_0}, using \eqref{eq:kappa3} and \eqref{eq:derivative1}, we see that, if $w=0$, then
$$\irn|\vp w_k|^p = o(1)\qquad\text{for every \ }\vp\in\cC^\infty_c(B_1(z)), \ z\in\rn.$$
Since $B_{C_1}(0)\subset D_k$ for sufficiently large $k$, arguing as in Case (1) we  obtain
$$\int_{B_{C_1}(0)}|w_k|^p=o(1).$$
This contradicts \eqref{eq:kappa3} and shows that $w\neq 0$.

If $\vp\in\cC^\infty_c(\rn)$, then $\supp(\vp)\subset D_k$ for all large enough $k$, and setting $\vp_k(z):=\vp\big(\frac{x-\xi_k}{\eps_k}\big)$ we get that
\begin{align*}
\irn\nabla\what w_k\cdot\nabla\vp-a_n\irn|\what w_k|^{p-2}w_k\vp &=\irn\nabla w_k\cdot\nabla\vp-a_n\irn|w_k|^{p-2}w_k\vp \\
& = J_b'(u_k)[\vp_k]=o(1).
\end{align*}
Passing to the limit yields $J_\infty'(w)=0$ in $D^{1,2}(\rn)'$. Hence, $w$ is a nontrivial $\phi$-equivariant solution to \eqref{problem_rn}. 

Fix a radial cut-off function $\chi\in\cC_c^\infty(\rn)$ such that $\chi(z)=1$ if $|z|\leq 1$ and $\chi(z)=0$ if $|z|\geq 2$ and set
$$\chi_k(z):=\chi\Big(\frac{\eps_k z}{r_k}\Big)\qquad\text{where \ }r_k:=\frac{1}{2}\dist(\xi_k,\partial\D).$$
Note that $w\chi_k\in D^{1,2}_0(D_k)$ and
\begin{align*}
\irn|\nabla(w(\chi_k-1)|^2 &\leq C\left(\int_{|z|\geq\frac{r_k}{\eps_k}}|\nabla w|^2 + \Big(\frac{\eps_k}{r_k}\Big)^2\int_{\frac{r_k}{\eps_k}\leq|z|\leq\frac{2r_k}{\eps_k}}w^2\right) \\
&\leq C\left(\int_{|z|\geq\frac{r_k}{\eps_k}}|\nabla w|^2 + \Big(\int_{|z|\geq\frac{r_k}{\eps_k}}w^p\Big)^\frac{2}{p}\right) = o(1).
\end{align*}
Therefore,
\begin{align*}
\int_{D_k}|\nabla(w_k-w\chi_k)|^2 &=\int_{D_k}|\nabla w_k|^2 - 2\int_{D_k}\nabla w_k\cdot\nabla(w\chi_k) + \int_{D_k}|\nabla(w\chi_k)|^2  \\
&=\int_{D_k}|\nabla w_k|^2 - 2\irn\nabla \what w_k\cdot\nabla(w\chi_k) + \int_{D_k}|\nabla(w\chi_k)|^2  \\
&=\int_{D_k}|\nabla w_k|^2 - \irn|\nabla w|^2 + o(1).
\end{align*}
It is also easy to see that
$$\int_{D_k}|w_k-w\chi_k|^p = \int_{D_k}|w_k|^p - \irn|w|^p + o(1).$$
As a consequence,
\begin{align}\label{eq:J(v)}
J_b(u_k)&=\frac{1}{2}\int_{D_k}|\nabla w_k|^2-\frac{a_n}{p}\int_{D_k}|w_k|^p - \frac{b}{q}\int_{\partial\D}|u_k|^q \nonumber \\
&=\frac{1}{2}\int_{D_k}|\nabla(w_k-w\chi_k)|^2-\frac{a_n}{p}\int_{D_k}|\nabla(w_k-w\chi_k)|^p - \frac{b}{q}\int_{\partial\D}|v_k|^q \nonumber\\
&\qquad+ \frac{1}{2}\irn|\nabla w|^2-\frac{a_n}{p}\irn|w|^p + o(1) \nonumber\\
&=J_b(v_k) + J_\infty(w) + o(1),
\end{align}
where
$$v_k(x):=u_k(x)-\eps_k^\frac{2-n}{2}w\Big(\frac{x-\xi_k}{\eps_k}\Big)\chi\Big(\frac{x-\xi_k}{r_k}\Big).$$
Similarly,
\begin{equation}\label{eq:J'(v)}
o(1)=J_b'(u_k)u_k=J_b'(v_k)v_k + J_\infty'(w)w + o(1)=J_b'(v_k)v_k + o(1).
\end{equation}
Note that $v_k$ is $\phi$-equivariant. It follows from \eqref{eq:J'(v)} that $v_k\to 0$ strongly in $D^{1,2}(\hn)$, as, otherwise, $J_b(v_k)+o(1)\geq\mu_b^\phi(\D)$, contradicting \eqref{eq:J(v)}. Therefore, 
$$\lim_{k\to\infty}\left\|u_k-\eps_k^\frac{2-n}{2}w\Big(\frac{\cdot - \xi_k}{\eps_k}\Big)\right\| = \lim_{k\to\infty}\|v_k\| = 0$$
and $\mu^\phi_b(\D)=J_\infty(w)\geq \mu^\phi_\infty$. Hence, in this case, we obtain statement $(II)$.

\underline{Case $(3)$}: $\xi_k=(0,1)$ for all $k\in\n$.

In this case, $\bigcup_{k\geq 1}D_k=\mathbb{H}:=\{(x',t)\in\r^{n-1}\times\r:t<1\}$ and the functions 
$$\what w_k(z):=\eps_k^\frac{n-2}{2}u_k^*(\eps_kz+\xi_k)$$
belong to $D^{1,2}_0(\mathbb{H})$. A subsequence satisfies $\what w_k\rh w$ weakly in $D^{1,2}(\rn)$, $\what w_k\to w$ a.e. in $\rn$, $\what w_k\to w$ in $L^2_{loc}(\rn)$ and $w\in D^{1,2}_0(\mathbb{H})$. If $\vp\in\cC^\infty_c(B_1(z))$ for some $z\in\rn$, then $\supp(\vp)\subset \eps_k^{-1}(\hn-\xi_k)=\{(x',t)\in\r^{n-1}\times\r:t>-\frac{1}{\eps_k}\} $ for large enough $k$ and, arguing as in Case (2), we show that $w\neq 0$.

On the other hand, if $\vp\in\cC_c^\infty(\mathbb{H})$, then $\supp(\vp)\subset D_k$ for large enough $k$, and setting $\vp_k(z):=\vp\big(\frac{x-\xi_k}{\eps_k}\big)$ we obtain
\begin{equation*}
\int_\mathbb{H}\nabla w_k\cdot\nabla\vp-a_n\int_\mathbb{H}|w_k|^{p-2}w_k\vp = J_b'(u_k)[\vp_k]=o(1).
\end{equation*}
Passing to the limit shows that $w$ is a nontrivial solution to the problem 
\begin{equation*}
\begin{cases}
 -\Delta u = a_n|u|^\frac{4}{n-2}u,\\
 u\in D^{1,2}_0(\mathbb{H}),
\end{cases}
\end{equation*}
in the half-space $\mathbb{H}$. It is well known that this problem does not have a nontrivial solution. Therefore, Case (3) cannot occur.

This completes the proof.
\end{proof}

\end{document}
